% Generated draft: requires mathematical and source review.
\documentclass[11pt]{article}
\usepackage[margin=25mm]{geometry}
\usepackage{fontspec}
\setmainfont{Latin Modern Roman}
\usepackage{amsmath,amssymb,mathtools,unicode-math}
\setmathfont{Latin Modern Math}
\usepackage{enumitem,needspace,xurl,hyperref}
\setlength{\parindent}{0pt}
\setlength{\parskip}{5pt}
\setlength{\emergencystretch}{3em}
\newcommand{\N}{\mathbb{N}}
\newcommand{\Z}{\mathbb{Z}}
\newcommand{\Q}{\mathbb{Q}}
\newcommand{\R}{\mathbb{R}}
\newcommand{\C}{\mathbb{C}}
\newcommand{\F}{\mathbb{F}}
\newcommand{\A}{\mathbb{A}}
\newcommand{\PP}{\mathbb{P}}
\newcommand{\E}{\mathbb{E}}
\newcommand{\eps}{\varepsilon}
\newcommand{\dd}{\,\mathrm{d}}
\begin{document}
\section*{UnsolvedMath ID: 30005213}
\textbf{Source problem code:} catalogue entry
\subsection{Definitions and mathematical statement}
Let \(\Omega\subset\R^2\) be open, the planar setting implicit in the supplied distributional formulation, and let \(S^1=\{v\in\R^2:|v|=1\}\). No boundary conditions, boundedness, or connectedness assumptions are supplied. A measurable field \(m:\Omega\to S^1\) satisfies the divergence-free eikonal constraint if
\[
|m|=1\quad\text{almost everywhere in }\Omega,
\qquad
\int_\Omega m\cdot\nabla\varphi\,\dd x=0
\quad\text{for every }\varphi\in C_c^\infty(\Omega).
\]
For \(3\le q<\infty\), write \(m\in B^{1/3}_{q,\mathrm{loc},\infty}(\Omega;\R^2)\) when \(m\in L^q_{\mathrm{loc}}\) and, for every open \(U\Subset\Omega\), there exists \(\delta_U>0\) with \(U+B_{\delta_U}(0)\subset\Omega\) such that
\[
\sup_{0<|h|<\delta_U}|h|^{-1/3}
\|m(\cdot+h)-m(\cdot)\|_{L^q(U)}<\infty.
\]

\textbf{Threshold question.} Determine whether, among divergence-free unit fields, the critical integrability index at Besov smoothness \(1/3\) is exactly \(q=3\): does
\[
\text{for every }3<q<\infty,\qquad
m\in B^{1/3}_{q,\mathrm{loc},\infty}(\Omega;\R^2),
\quad \operatorname{div}m=0,
\quad |m|=1\ \text{a.e.}
\quad\Longrightarrow\quad
m\text{ is rigid}
\]
hold, while the endpoint \(B^{1/3}_{3,\mathrm{loc},\infty}\) permits line singularities? Here ``rigid'' is the terminology of the supplied report extract. That extract describes the expected conclusion as allowing only a discrete set of point singularities, but does not define the singular set, specify regularity away from it, or give a complete definition of rigidity. Consequently this implication records the intended target, rather than a fully specified theorem. In particular, the endpoint is proposed as a dividing threshold; the supplied discussion does not assert that membership at the endpoint itself guarantees rigidity.

\textbf{Entropy formulation in the source.} With \(e(t)=(\cos t,\sin t)\), the extracted entropy condition is
\[
\frac{\dd}{\dd t}\Phi(e(t))\cdot e(t)=0
\qquad\text{for every }t\in\R.
\]
The source denotes the corresponding class of maps \(\Phi:S^1\to\R^2\) by \(\mathrm{ENT}\); its admissible smoothness class is not recoverable from the supplied text. It presents a conjecture with hypothesis
\[
\operatorname{div}(\Phi\circ m)\in L^r_{\mathrm{loc}}(\Omega)
\qquad\text{for every }\Phi\in\mathrm{ENT},
\]
and conclusion that \(m\) is rigid. The exponent is renamed \(r\) here to distinguish it from the Besov integrability index \(q\). The conjecture's range of \(r\) is missing from the extraction, so it cannot be supplied as an exact quantified assertion.

The extract attributes to its reference [7] an equivalence between this entropy hypothesis and \(m\in B^{1/3}_{3r,\mathrm{loc},\infty}\), for a range apparently bounded by \(1\) and \(4/3\); the endpoint conventions are corrupted. It attributes to its reference [2] an equivalence between the condition \(\operatorname{div}(\Phi\circ m)\in\mathcal M_{\mathrm{loc}}(\Omega)\) for every entropy and membership in \(B^{1/3}_{3,\mathrm{loc},\infty}\), where \(\mathcal M_{\mathrm{loc}}\) denotes locally finite signed Radon measures. These are supplied contextual claims whose precise hypotheses and bibliographic identities require review. The source further states that the endpoint class includes all finite limits of Aviles--Giga energy and can contain line singularities; the energy, convergence, and singularity conventions are not specified. It describes the rigidity conjecture as a weak consequence of the Aviles--Giga conjecture.
\subsection{Short English statement}
The intended question is whether every divergence-free unit vector field in \(B^{1/3}_{q,\mathrm{loc},\infty}\), for every \(q>3\), has the rigidity described in the source as permitting only discrete point singularities, although the endpoint \(q=3\) can permit line singularities. A complete theorem statement still requires the source's definition of rigidity and the missing entropy hypotheses. The supplied partially solved status and reported disk results are dated curation, not independently verified conclusions.
\subsection{Sources}
\begin{itemize}
\item Supplied local UnsolvedMath registry record, permanent ID \(30005213\), title \emph{Threshold Regularity for the Eikonal Equation}, exact source problem code \texttt{catalogue entry}, in the Oberwolfach Reports Open Problems collection. The record gives proposed year 2022 and no named proposer. Its labels \texttt{corrected\_verified} and \texttt{partially\_solved} are supplied registry metadata, not independent mathematical verification.
\item Supplied report citation: \emph{Calculus of Variations (2023)}, Oberwolfach Reports, DOI 10.4171/owr/2022/37, pp.\ 2129--2193. \url{https://doi.org/10.4171/owr/2022/37}. The extract identifies the relevant discussion as Problem 3, ``Threshold regularity of the eikonal equation.'' Its numbered references [2] and [7] are not bibliographically identified in the supplied record.
\item Supplied literature source: Lamy--Lorent--Peng, ``eikonal regularization.'' \url{https://arxiv.org/abs/2504.20933}.
\item Supplied literature source: Lorent--Peng, ``entropy production.'' \url{https://arxiv.org/abs/2104.01467}.
\item Supplied literature source: ``original Oberwolfach report.'' \url{https://doi.org/10.4171/owr/2022/37}.
\item Literature assessment dated exactly \texttt{2026-08-22}: the supplied curation labels the problem partially solved and states that the Aviles--Giga threshold-rigidity conjecture remains open for general \(p>3\), while a 2025 result proves the predicted regularizing effect at \(p=6\) and slightly below six in a disk setting. This is supplied dated curation derived from the OWR record, not an independently verified literature assessment. The exact exponent range, disk hypotheses, conclusion, and meaning of this \(p\) require checking against the cited source.
\end{itemize}
\end{document}
